\documentclass[reqno]{amsart}%
\usepackage{amssymb,mathtools,enumerate}
%\linenumbers,lineno

\usepackage{ifpdf}
\ifpdf
 \usepackage[hyperindex]{hyperref}%,pagebackref
\else
 \expandafter\ifx\csname dvipdfm\endcsname\relax
 \usepackage[hypertex,hyperindex]{hyperref}%
 \else
 \usepackage[dvipdfm,hyperindex]{hyperref}%
 \fi
\fi

\allowdisplaybreaks[4]
%\numberwithin{equation}{section}

\theoremstyle{plain}
\newtheorem{thm}{Theorem}
\newtheorem{cor}{Corollary}
\newtheorem{conj}{Conjecture}
\newtheorem{open}{Problem}
\newtheorem{lem}{Lemma}

\theoremstyle{remark}
\newtheorem{rem}{Remark}

\DeclareMathOperator{\td}{d\!}
\DeclareMathOperator{\te}{e}

%\usepackage{refcheck,cmdtrack}
%\overfullrule=5pt

\begin{document}

\title[Sectional curvature of Fisher--Rao manifold]
{The decrease in sectional curvature of the Fisher--Rao manifold of beta distributions}

\author[F. Qi]{Feng Qi}
\address{School of Mathematics and Physics, Hulunbuir University, Hulunbuir 021008, Inner Mongolia, China;
17709 Sabal Court, University Village, Dallas, TX 75252-8024, USA}
\email{\href{mailto: F. Qi <qifeng618@gmail.com>}{qifeng618@gmail.com}}
\urladdr{\url{https://orcid.org/0000-0001-6239-2968}}

\begin{abstract}
In the paper, the author develops a general convexity lemma and applies it to demonstrate the monotonic decrease of a function involving polygamma functions, arising from the sectional curvature of the Fisher--Rao manifold of beta distributions.
\end{abstract}

\keywords{sectional curvature, decrease, lower bound, Fisher--Rao manifold, beta distribution, polygamma function}

\subjclass[2010]{Primary 53C25; Secondary 26A48, 26A51, 33B15, 60E05}

\thanks{This paper was typeset using \AmS-\LaTeX}

\maketitle

\section{Introduction}
The Fisher--Rao manifold of beta distributions is the parameter space of the beta family endowed with the Fisher--Rao metric. This geometric structure arises in information geometry, a growing field that examines statistical models through the lens of differential geometry. The resulting manifold provides insight into the intrinsic geometry of beta distributions and supports deeper analysis of their statistical and inferential properties.
\par
Let $M=\{(x,y): x,y>0\}$ denote the first quadrant on $\mathbb{R}^2$. Let
\begin{equation*}
\td s^2=\psi'(x)\td x^2+\psi'(y)\td y^2-\psi'(x+y)(\td x+\td y)^2
\end{equation*}
be the Fisher metric and $M$ be equipped with $\td s^2$, where, for $\Re(z)>0$,
\begin{equation*}
\psi(z)=[\ln\Gamma(z)]'=\frac{\Gamma'(z)}{\Gamma(z)}\quad\text{and}\quad
\Gamma(z)=\int_0^{\infty}t^{z-1}\te^{-t}\td t.
\end{equation*}
\par
In~\cite[Proposition~3]{brigant2019fisherrao}, \cite[Proposition~14]{brigant2020fisherraoV1}, \cite[Proposition~15]{brigant2020fisherrao}, and~\cite[Proposition~15]{brigant2021fisherrao}, the authors computed the sectional curvature $K(x,y)$ of the Fisher metric $\td s^2$ on the beta manifold $\bigl(M,\td s^2\bigr)$ and obtained that
\begin{equation}\label{curvature-polygamma-exp}
K(x,y)=\frac{1}{4}\frac{\psi''(x)\psi''(y)\psi''(x+y) \Bigl[\frac{\psi'(x)}{\psi''(x)}+\frac{\psi'(y)}{\psi''(y)}-\frac{\psi'(x+y)}{\psi''(x+y)}\Bigr]}
{[\psi'(x)\psi'(x+y)+\psi'(y)\psi'(x+y)-\psi'(x)\psi'(y)]^2}, \quad x,y\in M.
\end{equation}
\par
In~\cite[Proposition~5]{brigant2019fisherrao}, \cite[Theorem~6]{brigant2020fisherraoV1}, \cite[Theorem~6]{brigant2020fisherrao}, and~\cite[Theorem~6]{brigant2021fisherrao}, the authors established that the sectional curvature $K(x,y)$ is negative and bounded from below, and they subsequently proposed the following conjecture.

\begin{conj}[{\cite[Proposition~5]{brigant2019fisherrao}, \cite[pp.~12--13]{brigant2020fisherraoV1}, \cite[p.~13--14]{brigant2020fisherrao}, and~\cite[p.~14]{brigant2021fisherrao}}]\label{Alice-lower-bound-conj}
For $x,y>0$, the sectional curvature $K(x,y)$ given in~\eqref{curvature-polygamma-exp} satisfies that
\begin{enumerate}
\item
it has a lower bound $-\frac{1}{2}$, accurately, $K(x,y)>-\frac{1}{2}$;
\item
it is decreasing in both $x$ and $y$.
\end{enumerate}
\end{conj}

In the review article~\cite{manifold-gamma.tex}, together with about ten papers cited therein, several functions involving the polygamma functions $\psi^{(n)}(x)$---originating from the study of the sectional curvature $K(x,y)$---were investigated extensively and deeply.
\par
Taking $y=x$ in~\eqref{curvature-polygamma-exp} yields that
\begin{equation}\label{K(x-x)-Eq}
\mathcal{K}(x)=K(x,x)=\frac{1}{4}\frac{\psi''(x)}{[\psi'(x)]^2} \frac{2\psi'(x)\psi''(2x)-\psi'(2x)\psi''(x)}{[\psi'(x)-2\psi'(2x)]^2}, \quad x\in(0,\infty).
\end{equation}
In~\cite[Theorem~4.1]{RIMA-D-20-01130.tex}, the sharp double inequality
\begin{equation}\label{sectional-doub-ineq}
0<\frac{2\psi'(x)\psi''(2x)-\psi'(2x)\psi''(x)}{[\psi'(x)-2\psi'(2x)]^2}<2, \quad x\in(0,\infty)
\end{equation}
was established. Consequently, the double inequality
\begin{equation}\label{sect-curvat-ineq}
0>\mathcal{K}(x)>-\frac{1}{2}, \quad x\in(0,\infty)
\end{equation}
was derived in~\cite[Theorem~5.1]{RIMA-D-20-01130.tex}. This may be viewed as the initial and partial confirmation of Conjecture~\ref{Alice-lower-bound-conj}.
\par
In this paper, we aim to prove that the sectional curvature $\mathcal{K}(x)$ defined in~\eqref{K(x-x)-Eq} is decreasing for $x\in(0,\infty)$. This result can be regarded as the second, though partial, confirmation of Conjecture~\ref{Alice-lower-bound-conj}.

\section{Two lemmas}
In order to attain the aim of this paper, we need the following two lemmas.

\begin{lem}\label{lem:general}
Let $f(x)$ be thrice continuously differentiable and $H(x)=-\frac{f(x)}{f'(x)}$ on $(0,\infty)$. If
$$
f(x)>0,\quad f'(x)<0, \quad 0<H'(x)<1, \quad H''(x)\geq 0
$$
on $(0,\infty)$, then, for every given $a>0$, the function
$$
F_a(x)=\frac{1}{1-\frac{f(x+a)}{f(x)}}
$$
is strictly convex on $(0,\infty)$.
\end{lem}

\begin{proof}
Define $\Lambda(x)=\ln\frac{f(x)}{f(x+a)}$ for $x\in(0,\infty)$.
Since $\frac{f'(x)}{f(x)}=-\frac1{H(x)}$, we have
$$
\Lambda(x)=\ln f(x)-\ln f(x+a)
=\int_{x+a}^{x}\frac{f'(t)}{f(t)}\td t
=
\int_x^{x+a}\frac{\td t}{H(t)}>0.
$$
Consequently, it follows that
$$
\Lambda'(x)=\frac1{H(x+a)}-\frac1{H(x)}<0
$$
and
$$
\Lambda''(x)
=\frac{H'(x)}{H^2(x)}-\frac{H'(x+a)}{H^2(x+a)}.
$$
Therefore, we acquire
\begin{equation}\label{eq:lambda-bound}
\begin{aligned}
\frac{\Lambda''(x)}{[\Lambda'(x)]^2}
&=\frac{H'(x)H^2(x+a)-H'(x+a)H^2(x)}{[H(x+a)-H(x)]^2}\\
&\leq \frac{H'(x)[H^2(x+a)-H^2(x)]}{[H(x+a)-H(x)]^2}\\
&=H'(x)\frac{H(x+a)+H(x)}{H(x+a)-H(x)}\\
&=H'(x)\coth\frac{L(x)}{2}, \quad x\in(0,\infty),
\end{aligned}
\end{equation}
where
$$
L(x)=\ln\frac{H(x+a)}{H(x)}>0, \quad x\in(0,\infty).
$$
Moreover, it is easy to see that
$$
L(x)=\int_x^{x+a}\frac{H'(t)}{H(t)}\td t
\geq H'(x)\int_x^{x+a}\frac{\td t}{H(t)}
=H'(x)\Lambda(x).
$$
Hence, we arrive at
\begin{equation}\label{LL-ineq}
\frac{\Lambda(x)}{2}\leq\frac{L(x)}{2H'(x)}, \quad x\in(0,\infty).
\end{equation}
\par
Let
$$
G(z)=z\coth z>0, \quad z\in(0,\infty).
$$
Direct differentiation gives
$$
G'(z)=\frac{\sinh z\cosh z-z}{\sinh^2z}
$$
and
$$
(\sinh z\cosh z-z)'=2\sinh^2z>0.
$$
This implies that $G'(z)>0$ and then $G(z)$ is strictly increasing for $z\in(0,\infty)$.

Since $0<H'(x)<1$, it follows that
$$
\frac{L(x)}{2H'(x)}>\frac{L(x)}{2}.
$$
The strict increase of $G(z)$ results in
$$
\coth\frac{L(x)}{2H'(x)}>H'(x)\coth\frac{L(x)}{2}.
$$
Since $\coth z$ is strictly decreasing on $(0,\infty)$ and the inequality~\eqref{LL-ineq}, we obtain
$$
\coth\frac{\Lambda(x)}{2}
\geq
\coth\frac{L(x)}{2H'(x)}
>
H'(x)\coth\frac{L(x)}{2}.
$$
Combining this with~\eqref{eq:lambda-bound} leads to the inequality
\begin{equation}\label{eq:key-inequality}
\frac{\Lambda''(x)}{[\Lambda'(x)]^2}<\coth\frac{\Lambda(x)}{2}, \quad x\in(0,\infty).
\end{equation}
\par
Now assume that
$$
q(x)=\te^{-\Lambda(x)}=\frac{f(x+a)}{f(x)}, \quad x\in(0,\infty).
$$
Then $0<q(x)<1$ and
$$
F_a(x)=\frac1{1-q(x)}.
$$
Furthermore, we have
$$
q'(x)=-q(x)\Lambda'(x)\quad\text{and}\quad
q''(x)=q(x)\bigl([\Lambda'(x)]^2-\Lambda''(x)\bigr).
$$
Thus, the second derivative
$$
F_a''(x)=\frac{q''(x)}{[1-q(x)]^2}+\frac{2[q'(x)]^2}{[1-q(x)]^3}
=
\frac{q(x)}{[1-q(x)]^2}
\left[
\frac{1+q(x)}{1-q(x)}[\Lambda'(x)]^2-\Lambda''(x)
\right].
$$
Since
$$
\frac{1+q(x)}{1-q(x)}
=
\coth\frac{\Lambda(x)}{2},
$$
we obtain
$$
F_a''(x)
=
\frac{\te^{-\Lambda(x)}[\Lambda'(x)]^2}{[1-\te^{-\Lambda(x)}]^2}
\biggl[\coth\frac{\Lambda(x)}{2}-\frac{\Lambda''(x)}{[\Lambda'(x)]^2}\biggr].
$$
The bracket is strictly positive by~\eqref{eq:key-inequality}.
Therefore, it follows that $F_a''(x)>0$. The proof of Lemma~\ref{lem:general} is thus complete.
\end{proof}

Applying Lemma~\ref{lem:general} to $f(x)=\bigl|\psi^{(n)}(x)\bigr|$, we obtain the second lemma below.

\begin{lem}\label{OpenProb-convex}
For fixed $\alpha>0$ and $n\in\mathbb{N}$, the function
\begin{equation}\label{Phi_{n-alpha}(x)}
\Phi_{n,\alpha}(x)=\frac{1}{1-\frac{\psi^{(n)}(x+\alpha)}{\psi^{(n)}(x)}}
\end{equation}
is strictly increasing and strictly convex on $(0,\infty)$.
\end{lem}

\begin{proof}
In~\cite[p.~260, Entry~6.4.1]{abram}, the integral representation
\begin{equation}\label{psin}
\psi^{(n)}(z)=(-1)^{n+1}\int_{0}^{\infty}\frac{t^n}{1-\te^{-t}}\te^{-zt}\td t
\end{equation}
is listed for $\Re(z)>0$ and $n\in\mathbb{N}$.
Letting 
$$
f(x)=\bigl|\psi^{(n)}(x)\bigr|=(-1)^{n+1}\psi^{(n)}(x)>0
$$ 
for $x\in(0,\infty)$ in Lemma~\ref{lem:general} gives that
\begin{equation*}
f(x)=\int_{0}^{\infty}\frac{t^n}{1-\te^{-t}}\te^{-xt}\td t>0, \quad x\in(0,\infty).
\end{equation*}
Accordingly,
\begin{equation*}
f'(x)=-\int_{0}^{\infty}\frac{t^{n+1}}{1-\te^{-t}}\te^{-xt}\td t<0, \quad x\in(0,\infty)
\end{equation*}
and
\begin{equation*}
H(x)=\frac{\int_{0}^{\infty}\frac{t^n}{1-\te^{-t}}\te^{-xt}\td t} {\int_{0}^{\infty}\frac{t^{n+1}}{1-\te^{-t}}\te^{-xt}\td t}
=-\frac{\psi^{(n)}(x)}{\psi^{(n+1)}(x)}, \quad x\in(0,\infty).
\end{equation*}
\par
Differentiating yields
\begin{equation*}
H'(x)=\frac{\psi^{(n)}(x)\psi^{(n+2)}(x)}{[\psi^{(n+1)}(x)]^2}-1
\end{equation*}
and
\begin{equation*}
H''(x)=\frac{\psi^{(n+2)}(x)}{\psi^{(n+1)}(x)}+\frac{\psi^{(n)}(x) \psi^{(n+3)}(x)}{[\psi^{(n+1)}(x)]^2} -\frac{2\psi^{(n)}(x)[\psi^{(n+2)}(x)]^2}{[\psi^{(n+1)}(x)]^3}.
\end{equation*}
\par
The first statement in~\cite[Theorem~2]{Yang-JMAA-2017} reads that the function $\frac{[\psi^{(n+1)}(x)]^2}{\psi^{(n)}\psi^{(n+2)}(x)}$ for $n\in\mathbb{N}$ is strictly decreasing from $(0,\infty)$ onto $\bigl(\frac{n}{n+1},\frac{n+1}{n+2}\bigr)$.
See also~\cite[Corollary~2.3]{Alzer-Wells-SIAM-98}, \cite[Section~3.5, pp.~10--12]{Qi-Agar-Surv-JIA.tex}, and~\cite[p.~3, Eq.~(8)]{Qi-Nantomah-Lim.tex}. Consequently, the first derivative $H'(x)$ is strictly increasing from $(0,\infty)$ onto $\bigl(\frac{n+2}{n+1}-1, \frac{n+1}{n}-1\bigr)=\bigl(\frac{1}{n+1},\frac{1}{n}\bigr)$ and then the second derivative $H''(x)$ is positive. As a result, all the hypotheses of Lemma~\ref{lem:general} are satisfied and the function $\Phi_{n,\alpha}(x)$ defined by~\eqref{Phi_{n-alpha}(x)} is strictly convex on $(0,\infty)$.
\par
Straightforward computation yields
\begin{align*}
\Phi_{n,\alpha}'(x)&=\frac{\psi^{(n+1)}(x)\psi^{(n+1)}(x+\alpha)}{[\psi^{(n)}(x)-\psi^{(n)}(x+\alpha)]^2} \biggl[\frac{\psi^{(n)}(x)}{\psi^{(n+1)}(x)} -\frac{\psi^{(n)}(x+\alpha)}{\psi^{(n+1)}(x+\alpha)}\biggr]\\
&=\frac{x^n\bigl[x^{n+2}\psi^{(n+1)}(x)\bigr]\psi^{(n+1)}(x+\alpha)}{[x^{n+1}\psi^{(n)}(x)-x^{n+1}\psi^{(n)}(x+\alpha)]^2} \biggl[\frac{x[x^{n+1}\psi^{(n)}(x)]}{x^{n+2}\psi^{(n+1)}(x)} -\frac{\psi^{(n)}(x+\alpha)}{\psi^{(n+1)}(x+\alpha)}\biggr]\\
&\to0
\end{align*}
as $x\to0^+$, where we used the limit
\begin{equation*}%\label{polyG-lim-to0plus}
\lim_{x\to0^+}\bigl[x^k\psi^{(k-1)}(x)\bigr]=(-1)^{k}(k-1)!, \quad k\in\mathbb{N}
\end{equation*}
in~\cite[Lemma~2.1]{RIMA-D-20-01130.tex}. Consequently, from the facts that $\Phi_{n,\alpha}''(x)>0$ and $\Phi_{n,\alpha}'(x)$ is strictly increasing on $(0,\infty)$, it follows that $\Phi_{n,\alpha}'(x)>0$ on $(0,\infty)$, and then that the function $\Phi_{n,\alpha}(x)$ is strictly increasing on $(0,\infty)$.
The proof of Lemma~\ref{OpenProb-convex} is thus complete.
\end{proof}

\begin{rem}
From the proof of Lemma~\ref{OpenProb-convex}, we can conclude that, for $n\in\mathbb{N}$, the inequality
\begin{equation*}
\frac{1}{2}\biggl[\frac{\psi^{(n+2)}(x)}{\psi^{(n+1)}(x)}+\frac{\psi^{(n)}(x) \psi^{(n+3)}(x)}{[\psi^{(n+1)}(x)]^2}\biggr] >\frac{\psi^{(n)}(x)[\psi^{(n+2)}(x)]^2}{[\psi^{(n+1)}(x)]^3}
\end{equation*}
is valid on $(0,\infty)$.
\end{rem}

\begin{rem}
Lemma~\ref{OpenProb-convex} is a complete nice answer to the question announced at \url{https://mathoverflow.net/q/396837} (accessed on 5 September 2026).
\end{rem}

\begin{rem}
Lemma~4 in~\cite{Yang-Tian-JIA-2017} states that, if the ratio $\frac{A(t)}{B(t)}$ is increasing, then the ratio $\frac{\int_{0}^{\infty}A(t)\te^{-xt}\td t}{\int_{0}^{\infty}B(t)\te^{-xt}\td t}$ is decreasing on $(0,\infty)$, where the functions $A(t)$ and $B(t)\ne0$ are defined on $(0,\infty)$ such that their Laplace transforms exist. A more general ratio monotonicity rule was established in~\cite[Lemma~9 and Remark~15]{Alice-y-x-conj-one.tex} together with~\cite[Remark~7.2]{Ouimet-LCM-BKMS.tex}. By virtue of the integral representation~\eqref{psin}, we have
\begin{equation*}
\frac{\psi^{(n)}(x+\alpha)}{\psi^{(n)}(x)}=\frac{\int_{0}^{\infty}\frac{t^n\te^{-\alpha t}}{1-\te^{-t}}\te^{-xt}\td t} {\int_{0}^{\infty}\frac{t^n}{1-\te^{-t}}\te^{-xt}\td t}.
\end{equation*}
Accordingly, in light of the above-mentioned ratio monotonicity rule~\cite[Lemma~9]{Alice-y-x-conj-one.tex} or~\cite[Lemma~4]{Yang-Tian-JIA-2017}, we arrive at the increase of $\frac{\psi^{(n)}(x+\alpha)}{\psi^{(n)}(x)}$ in $x\in(0,\infty)$. As a result, the increase (without strictness) of the function $\Phi_{n,\alpha}(x)$ on $(0,\infty)$ is alternatively proved.
\end{rem}

\section{The decrease of $\mathcal{K}(x)$}
Now we are in a position to prove our main result.

\begin{thm}\label{sect-curvat-decr-thm}
The sectional curvature $\mathcal{K}(x)$ defined in~\eqref{K(x-x)-Eq} is strictly decreasing on $(0,\infty)$. Consequently, the inequality~\eqref{sect-curvat-ineq} is sharp.
\end{thm}

\begin{proof}
In~\cite[Theorem~3.2]{RIMA-D-20-01130.tex}, the second factor $\frac{\psi''(x)}{[\psi'(x)]^2}$ in~\eqref{K(x-x)-Eq} was proved to be decreasing from $(0,\infty)$ onto $(-1,0)$.
\par
The third factor in~\eqref{K(x-x)-Eq} can be reformulated as
\begin{multline*}
\frac{2\psi'(x)\psi''(2x)-\psi'(2x)\psi''(x)}{[\psi'(x)-2\psi'(2x)]^2}
=\frac{\psi'(x)\psi''\bigl(x+\frac{1}{2}\bigr) -\psi'\bigl(x+\frac{1}{2}\bigr)\psi''(x)} {\bigl[\psi'(x)-\psi'\bigl(x+\frac{1}{2}\bigr)\bigr]^2}\\
=\biggl[\frac{\psi'\bigl(x+\frac12\bigr)}{\psi'(x)}\biggr]' \frac{1}{\Bigl[1-\frac{\psi'(x+\frac{1}{2})}{\psi'(x)}\Bigr]^2}
=\left[\frac{1}{1-\frac{\psi'(x+\frac{1}{2})}{\psi'(x)}}\right]',
\end{multline*}
where we utilized the first two derivatives
\begin{equation*}%\label{digamma-2x-eq}
\psi'(2z)=\frac{1}{4}\biggl[\psi'(z)+\psi'\biggl(z+\frac{1}{2}\biggr)\biggr]
\end{equation*}
and
\begin{equation*}%\label{tetragamma-2x-eq}
\psi''(2z)=\frac{1}{8}\biggl[\psi''(z)+\psi''\biggl(z+\frac{1}{2}\biggr)\biggr]
\end{equation*}
of the duplication formula
\begin{equation*}
\psi(2z)=\frac{1}{2}\biggl[\psi(z)+\psi\biggl(z+\frac{1}{2}\biggr)\biggr]+\ln2
\end{equation*}
in~\cite[p.~259, 6.3.8]{abram}. In view of Lemma~\ref{OpenProb-convex} applied to $n=1$ and $\alpha=\frac{1}{2}$, we conclude that
\begin{equation*}
\biggl[\frac{2\psi'(x)\psi''(2x)-\psi'(2x)\psi''(x)}{[\psi'(x)-2\psi'(2x)]^2}\biggr]'
=\left[\frac{1}{1-\frac{\psi'(x+\frac{1}{2})}{\psi'(x)}}\right]''
>0
\end{equation*}
for $x\in(0,\infty)$. Accordingly, with the aid of the sharp inequality~\eqref{sectional-doub-ineq}, the third factor in~\eqref{K(x-x)-Eq} is positive and strictly increasing on $(0,\infty)$.
\par
In conclusion, the sectional curvature $\mathcal{K}(x)$ in~\eqref{K(x-x)-Eq} is strictly decreasing on $(0,\infty)$. Combining this with the limits
\begin{equation*}
\lim_{x\to0^+}\mathcal{K}(x)=0 \quad\text{and}\quad \lim_{x\to\infty}\mathcal{K}(x)=-\frac{1}{2}
\end{equation*}
in the proof of~\cite[Theorem~5.1]{RIMA-D-20-01130.tex} demonstrates the sharpness of the inequality~\eqref{sect-curvat-ineq}. The proof of Theorem~\ref{sect-curvat-decr-thm} is thus complete.
\end{proof}

\section{More remarks}
In this section, we give three more remarks on our main result.

\begin{rem}
In~\cite[Remark~6.5]{Alice-y-x-curv-notes.tex}, the author guessed that the function
\begin{equation*}
\frac{2\psi'(x)\psi''(2x)-\psi'(2x)\psi''(x)}{[\psi'(x)-2\psi'(2x)]^2}
\end{equation*}
should be increasing in $x>0$. This implies that the function
\begin{equation*}
\frac{1}{1-\frac{\psi'(x+\frac{1}{2})}{\psi'(x)}}=\frac{\psi'(x)}{\psi'(x)-\psi'(x+\frac{1}{2})}
\end{equation*}
should be convex in $x\in(0,\infty)$. Lemma~\ref{OpenProb-convex} and the proof of Theorem~\ref{sect-curvat-decr-thm} confirm this guess positively.
\end{rem}

\begin{rem}
From the double inequality
\begin{equation}\label{psi(2k)-doub-ineq}
-\frac{1}{2}\frac{(2k)!}{(k-1)!k!}<\frac{\psi^{(2k)}(x)}{[\psi^{(k)}(x)]^2}<0, \quad x\in(0,\infty)
\end{equation}
in~\cite[Theorem~3.2]{RIMA-D-20-01130.tex}, it follows that
\begin{equation}\label{psi(2k=2)-doub-ineq}
-1<\frac{\psi''(x)}{[\psi'(x)]^2}<0, \quad x\in(0,\infty).
\end{equation}
In~\cite[Theorem~11]{Alice-y-x-conj-one.tex}, the double inequalities in~\eqref{psi(2k)-doub-ineq} and~\eqref{psi(2k=2)-doub-ineq} were generalized to a decreasing monotonicity of the function
\begin{equation}\label{Q(m-n)(x)-Q(ijlm)(x)-dfn}
Q_{m,n}(x)=\frac{\psi^{(m+n)}(x)}{\psi^{(m)}(x)\psi^{(n)}(x)}, \quad m,n\ge1
\end{equation}
from $(0,\infty)$ onto $\bigl(-\frac{(m+n-1)!} {(m-1)!(n-1)!},0\bigr)$. See also~\cite[Sections~9 and~12]{manifold-gamma.tex}. In particular, taking $m=n=1$ in~\eqref{Q(m-n)(x)-Q(ijlm)(x)-dfn} deduces that, the function $\frac{\psi''(x)}{[\psi'(x)]^2}$, the second factor in~\eqref{K(x-x)-Eq}, is decreasing from $(0,\infty)$ onto $(-1,0)$.
\end{rem}

\begin{rem}
The sectional curvature $K(x,y)$ can be rearranged as
\begin{equation*}
K(x,y)=\frac{1}{4}\frac{\psi''(x)}{[\psi'(x)]^2}\frac{\psi''(y)}{[\psi'(y)]^2}\frac{\psi''(x+y)}{[\psi'(x+y)]^2} \frac{\frac{\psi'(x)}{\psi''(x)}+\frac{\psi'(y)}{\psi''(y)}-\frac{\psi'(x+y)}{\psi''(x+y)}}
{\bigl[\frac{1}{\psi'(x)}+\frac{1}{\psi'(y)}-\frac{1}{\psi'(x+y)}\bigr]^2}.
\end{equation*}
Theorem~3.1 in~\cite{RIMA-D-20-01130.tex} reads that the function $\frac{\psi''(x)}{[\psi'(x)]^2}$ is decreasing from $(0,\infty)$ onto $(-1,0)$. Accordingly, 
\begin{enumerate}
\item
in order to confirm the decreasing monotonicity in Conjecture~\ref{Alice-lower-bound-conj}, it suffices to show that the function
\begin{align*}
Q(x,y)&=\frac{\psi''(x)}{[\psi'(x)]^2}\frac{\frac{\psi'(x)}{\psi''(x)}+\frac{\psi'(y)}{\psi''(y)}-\frac{\psi'(x+y)}{\psi''(x+y)}}
{\Bigl[\frac{1}{\psi'(x)}+\frac{1}{\psi'(y)}-\frac{1}{\psi'(x+y)}\Bigr]^2}\\
&=\frac{\psi''(x)}{[\psi'(x)]^2}\frac{\frac{\psi'(x+y)}{\psi''(x+y)}}{\Bigl[\frac{1}{\psi'(x+y)}\Bigr]^2}\frac{\frac{\frac{\psi'(x)}{\psi''(x)}+\frac{\psi'(y)}{\psi''(y)}} {\frac{\psi'(x+y)}{\psi''(x+y)}}-1}
{\biggl[\frac{\frac{1}{\psi'(x)}+\frac{1}{\psi'(y)}}{\frac{1}{\psi'(x+y)}}-1\biggr]^2}
\end{align*}
is decreasing in $x\in(0,\infty)$ for fixed $y\in(0,\infty)$;
\item
in order to confirm the lower bound $-\frac{1}{2}$ in Conjecture~\ref{Alice-lower-bound-conj}, it suffices to show that the double inequality $0<Q(x,y)<2$ is valid and sharp on $(0,\infty)$.
\end{enumerate}
\end{rem}

\section{Declarations}

\paragraph{\bf Authors' Contributions}
All authors contributed equally to the manuscript and read and approved the final manuscript. 

\paragraph{\bf Funding}
The author was partially supported by the Natural Science Foundation of Inner Mongolia Autonomous Region (Grant No.~2025QN01041) and by the Youth Project of Hulunbuir City for Basic Research and Applied Basic Research (Grant No.~GH2024020).

\paragraph{\bf Acknowledgements}
The author gratefully acknowledges Yun-Peng Meng, Ph.D. student in the Department of Mathematics at Texas A\&M University, for his helful assistance in the proof of Lemma~\ref{lem:general} of this paper.

\paragraph{\bf Institutional Review Board Statement}
Not applicable. 

\paragraph{\bf Informed Consent Statement}
Not applicable. 

\paragraph{\bf Ethical Approval}
The conducted research is not related to either human or animal use. 

\paragraph{\bf Availability of Data and Material}
Data sharing is not applicable to this article as no new data were created or analyzed in this study. 

\paragraph{\bf Competing Interests}
The authors declare that they have no conflict of competing interests. 

\paragraph{\bf Use of AI Tools Declaration}
The authors declare they have not used Artificial Intelligence (AI) tools in the creation of this article. 

\begin{thebibliography}{99}

\bibitem{abram}
M. Abramowitz and I. A. Stegun (Eds), \textit{Handbook of Mathematical Functions with Formulas, Graphs, and Mathematical Tables}, National Bureau of Standards, Applied Mathematics Series \textbf{55}, Reprint of the 1972 edition, Dover Publications, Inc., New York, 1992.

\bibitem{Alzer-Wells-SIAM-98}
H. Alzer and J. Wells, \emph{Inequalities for the polygamma functions}, SIAM J. Math. Anal. \textbf{29} (1998), no.~6, 1459\nobreakdash--1466. DOI: \url{https://doi.org/10.1137/S0036141097325071}.

\bibitem{brigant2019fisherrao}
A. Le Brigant and S. Puechmorel, \textit{The Fisher-Rao geometry of beta distributions applied to the study of canonical moments}, arXiv:1904.08247. URL: \url{https://arxiv.org/abs/1904.08247}. DOI: \url{https://doi.org/10.48550/arXiv.1904.08247}.

\bibitem{brigant2020fisherraoV1}
A. Le Brigant, S. C. Preston, and S. Puechmorel, \emph{Fisher--Rao geometry of Dirichlet distributions}, arXiv:2005.05608v1. URL: \url{https://arxiv.org/abs/2005.05608v1}. DOI: \url{https://doi.org/10.48550/arXiv.2005.05608}.

\bibitem{brigant2020fisherrao}
A. Le Brigant, S. C. Preston, and S. Puechmorel, \emph{Fisher--Rao geometry of Dirichlet distributions}, arXiv:2005.05608v2. URL: \url{https://arxiv.org/abs/2005.05608v2}. DOI: \url{https://doi.org/10.48550/arXiv.2005.05608}.

\bibitem{brigant2021fisherrao}
A. Le Brigant, S. C. Preston, and S. Puechmorel, \emph{Fisher--Rao geometry of Dirichlet distributions}, Differential Geom. Appl. \textbf{74} (2021), Paper No.~101702, 16~pp. DOI: \url{https://doi.org/10.1016/j.difgeo.2020.101702}.

\bibitem{Alice-y-x-curv-notes.tex}
F. Qi, \textit{Alternative proofs for monotonicity of some functions related to sectional curvature of Fisher--Rao manifold of beta distributions}, in: H. Dutta, N. Ahmed, R. P. Agarwal (eds), \emph{Applied Nonlinear Analysis and Soft Computing}, ANASC 2020; \emph{Advances in Intelligent Systems and Computing}, Vol.~1437, Chapter~23, pp.~401\nobreakdash--418; Springer, Singapore, 2023; available online at \url{https://doi.org/10.1007/978-981-19-8054-1_23}.

\bibitem{Alice-y-x-conj-one.tex}
F. Qi, \emph{Decreasing properties of two ratios defined by three and four polygamma functions}, C. R. Math. Acad. Sci. Paris \textbf{360} (2022), 89\nobreakdash--101. DOI: \url{https://doi.org/10.5802/crmath.296}.

\bibitem{RIMA-D-20-01130.tex}
F. Qi, \textit{Lower bound of sectional curvature of Fisher--Rao manifold of beta distributions and complete monotonicity of functions involving polygamma functions}, Results Math. \textbf{76} (2021), no.~4, Paper No.~217, 16~pp. DOI: \url{https://doi.org/10.1007/s00025-021-01530-2}.

\bibitem{Qi-Agar-Surv-JIA.tex}
F. Qi and R. P. Agarwal, \textit{On complete monotonicity for several classes of functions related to ratios of gamma functions}, J. Inequal. Appl. \textbf{2019}, Paper No.~36, 42~pp. DOI: \url{https://doi.org/10.1186/s13660-019-1976-z}.

\bibitem{manifold-gamma.tex}
F. Qi and R. P. Agarwal, \textit{Several functions originating from Fisher--Rao geometry of Dirichlet distributions and involving polygamma functions}, Mathematics \textbf{12} (2024), no.~1, Art.~44, 21~pp. DOI: \url{https://doi.org/10.3390/math12010044}.

\bibitem{Ouimet-LCM-BKMS.tex}
F. Qi, W.-H. Li, S.-B. Yu, X.-Y. Du, and B.-N. Guo, \textit{A ratio of finitely many gamma functions and its properties with applications}, Rev. R. Acad. Cienc. Exactas F\'{\i}s. Nat. Ser. A Mat. RACSAM \textbf{115} (2021), no.~2, Paper No.~39, 14~pp. DOI: \url{https://doi.org/10.1007/s13398-020-00988-z}.

\bibitem{Qi-Nantomah-Lim.tex}
F. Qi, D. Lim, and K. Nantomah, \emph{Monotonicity and positivity of several functions involving ratios and products of polygamma functions}, J. Inequal. Appl. \textbf{2025}, Paper No.~5, 10~pp. DOI: \url{https://doi.org/10.1186/s13660-024-03245-8}.

\bibitem{Yang-JMAA-2017}
Z.-H. Yang, \emph{Some properties of the divided difference of psi and polygamma functions}, J. Math. Anal. Appl. \textbf{455} (2017), no.~1, 761\nobreakdash--777. DOI: \url{https://doi.org/10.1016/j.jmaa.2017.05.081}.

\bibitem{Yang-Tian-JIA-2017}
Z.-H. Yang and J.-F. Tian, \emph{Monotonicity and inequalities for the gamma function}, J. Inequal. Appl. \textbf{2017}, Paper No.~317, 15~pp. DOI: \url{https://doi.org/10.1186/s13660-017-1591-9}.

\end{thebibliography}

\end{document} 